\documentclass[a4paper]{article}
\usepackage{graphicx}
\usepackage{amsmath,amssymb}
\usepackage{hyperref}
\usepackage{minted}
\usepackage{multirow}

\title{self Awareness}
\date{}
\begin{document}
    \maketitle
    see active inference on markov blanket that helps with drawing a boundry between self and environment
    
    In the philosophy of self, self-awareness is the awareness and reflection of one's own personality or individuality, including traits, feelings, and behaviors. \url{https://en.wikipedia.org/wiki/Self-awareness}
    
    \section{From the neuroscience point of view}
    Modern neuroscience treats self-awareness not as the product of a single ``center,`` but as the emergent behavior of interacting brain systems. Functional MRI, lesion, and connectivity studies implicate a distributed network—particularly the medial prefrontal cortex, anterior cingulate cortex, posterior cingulate cortex, and temporoparietal junction—in processes of self-reflection and self-modelling. These regions overlap substantially with the brain's default mode network and engage in metacognitive monitoring—feedback loops in which predictions about internal and external states are continuously compared with incoming sensory and emotional input. Such neural machinery underlies the experience of being aware that one is aware.\cite{fleming-frith-2014-metacognition}

    Experimental evidence suggests that self-awareness depends on the brain's capacity for metacognition—the monitoring and evaluation of its own processes. Such ``monitoring systems`` continually compare predicted sensory and emotional states with actual input, generating the experience of being aware of awareness itself. This feedback architecture allows the brain to notice discrepancies between expectation and perception, forming the foundation of conscious self-reflection. This recursive feedback process gives rise to the sensation of being aware of awareness, sometimes described as a "mirror of mirrors" within consciousness.\cite{garcia-2013-stuck-in-the-moment}


    
    It is the ability of the agent of itself and its distinction from the world
    
    compare with self conciouness
    \begin{minted}{python}
        class Awareness:
            def __init__(self):
                pass
            
            def discovered_actuators(self) -> list:
                pass
            
            def discovered_sensors(self) -> list:
                pass
            
            def what_am_i_doing_now(self) -> str:
                pass
    \end{minted}
    
    
    \bibliographystyle{plain}
    \bibliography{/home/donkarlo/Dropbox/repo/shared_research/bibliography/bibliography.bib}
\end{document}